\section{Introduction}
\label{sec:intro}

Background music fundamentally shapes the emotional experience of film, games,
podcasts, and audiobooks.
Yet composing scene-appropriate music requires significant musical expertise,
and manually selecting tracks for every new scene is prohibitively slow.
Automating this process---\emph{generating} music that is semantically
consistent with a given visual or textual scene---is therefore both practically
valuable and scientifically challenging.

Recent advances in text-to-audio diffusion~\cite{evans2024stable,
liu2023audioldm2, copet2023simple} have made it possible to synthesise
high-quality music from a natural-language prompt.
The remaining bottleneck is \emph{semantic bridging}: translating an
unstructured scene description (or image) into the rich, musically-specific
language these models require.
Existing approaches either rely on large language models (LLMs) at inference
time~\cite{liu2023audioldm2, agostinelli2023musiclm}, which is slow and
resource-intensive, or train end-to-end models that collapse the intermediate
representation, offering little interpretability or user control.

We propose a different strategy: \textbf{knowledge distillation through a
structured intermediate representation}.
A 7B-parameter teacher LLM (Qwen2.5~\cite{qwen2025qwen25}) labels a large
procedurally-generated scene corpus with structured \emph{MusicDescriptor}
JSON objects \emph{offline}.
A compact student MLP is then trained to predict these descriptors directly
from frozen CLAP~\cite{elizalde2023clap} embeddings.
At inference, no LLM is needed: the student runs in milliseconds, replacing
billions of parameters with fewer than one million.

The \emph{MusicDescriptor} is the linchpin of our approach.
It encodes 12 musically-grounded attributes---mood, energy, valence, tempo,
key mode, harmonic tension, texture density, instrumentation, rhythm style,
structure, production style, and dynamics profile---that a \texttt{prompt()}
method serialises into six natural-language sentences, conditioning
Stable Audio Open~\cite{evans2024stable}.
A final CLAP-based re-ranking step selects the best among $N$ generated
candidates.

\noindent\textbf{Contributions.}
\begin{enumerate}
  \item A structured \emph{MusicDescriptor} schema bridging multimodal scene
        understanding and text-to-audio synthesis.
  \item An efficient two-phase knowledge-distillation framework that replaces
        7B-parameter LLM inference with a sub-million-parameter MLP.
  \item A CLAP-based self-supervised re-ranking mechanism.
  \item Three demonstrated applications: image-to-ambient-music, adaptive
        audiobook narration, and image-to-video with a synchronised soundtrack.
  \item Quantitative benchmarks showing the distilled model (CLAP sim.\ 0.38)
        exceeds direct LLM inference (0.35) while eliminating the LLM from the
        inference graph.
\end{enumerate}

